\chapter{Unix 的工具传统}
\chapterlead{命令行工具的形态由一组四十年前形成的约定决定：单个工具只做一件事，工具之间用文本流连接，成功时不输出，配置放在固定位置。本章先用一条管道演示这些约定的作用，再说明它们的来源与维护方式。}

\section{起步：一条管道的构成}

在终端中执行下列一行。

\begin{shellsession}
$ history | cut -d' ' -f2- | sort | uniq -c | sort -rn | head -5
   87 ls
   61 cd
   42 vim
   29 git
   18 systemctl
\end{shellsession}

输出左侧是出现次数，右侧是输入过的命令。这一行由六个程序依次处理同一批数据，每个程序只承担一个步骤。

\begin{itemize}
  \item \texttt{history}：输出命令历史，一行一条，行首带序号。
  \item \texttt{cut -d' ' -f2-}：以空格为分隔符，从第二段取到行尾，去掉序号。\texttt{-d} 后引号内的一个空格即为分隔符；省略它，整条命令的结果会完全不同。
  \item \texttt{sort}：按字典序排列，使相同的行相邻。
  \item \texttt{uniq -c}：合并相邻的重复行，并在行首写入出现次数。
  \item \texttt{sort -rn}：按数值从大到小重新排序。
  \item \texttt{head -5}：只输出前五行。
\end{itemize}

竖线称为管道（pipe），作用是把左侧程序的标准输出接到右侧程序的标准输入。系统中没有名为“统计常用命令”的程序；结果是六个通用工具在一条命令中组合出来的。

\begin{figure}[H]
\centering
\begin{tikzpicture}[font=\normalsize,
  cmd/.style={draw=teal!60, fill=teal!7, rounded corners=1.5pt,
              inner sep=6pt, text=teal!85, font=\ttfamily\normalsize, anchor=west},
  dat/.style={draw=rule, fill=paper, rounded corners=1.5pt,
              inner sep=6pt, font=\ttfamily\normalsize, text=ink!90, anchor=west},
  arr/.style={-{Latex[length=1.6mm]}, teal, thick},
  note/.style={font=\small, text=muted, anchor=west}]

  \node[cmd] (c1) at (0,0)      {history};
  \node[dat] (d1) at (62mm,0)   {\ \ 512\ \ ls -la};
  \node[note] at (62mm,-7mm) {每行一条命令，行首带序号};

  \node[cmd] (c2) at (0,-22mm)  {cut -d' ' -f2-};
  \node[dat] (d2) at (62mm,-22mm) {ls -la};
  \node[note] at (62mm,-29mm) {按空格切段，取第二段到行尾};

  \node[cmd] (c3) at (0,-44mm)  {sort};
  \node[dat] (d3) at (62mm,-44mm) {cd / vim / ls -la / ls};
  \node[note] at (62mm,-51mm) {按字典序排列，相同的行相邻};

  \node[cmd] (c4) at (0,-66mm)  {uniq -c};
  \node[dat] (d4) at (62mm,-66mm) {\ \ 61 cd\ \ \ \ 87 ls};
  \node[note] at (62mm,-73mm) {相邻重复行合并，行首加计数};

  \node[cmd] (c5) at (0,-88mm)  {sort -rn};
  \node[dat] (d5) at (62mm,-88mm) {\ \ 87 ls\ \ \ \ 61 cd};
  \node[note] at (62mm,-95mm) {按数值从大到小重排};

  \node[cmd] (c6) at (0,-110mm)  {head -5};
  \node[dat] (d6) at (62mm,-110mm) {最常用的五条命令};
  \node[note] at (62mm,-117mm) {截断输出};

  \foreach \a/\b in {c1/c2,c2/c3,c3/c4,c4/c5,c5/c6} {
    \draw[arr] (\a.south) -- (\b.north);
  }
\end{tikzpicture}
\caption{一条管道中的六个步骤。左列是程序，右列是该步骤看到的数据}
\end{figure}

\section{一切皆为文件}

Unix 的一项设计决策是把“文件”作为统一的读写对象。

\begin{itemize}
  \item 磁盘上的文档是文件。
  \item 键盘与屏幕在 \texttt{/dev} 下有对应的设备文件。
  \item 运行中进程的信息在 \texttt{/proc} 下，一个进程一个目录。
  \item 内核参数、硬件状态、随机数都挂在某个路径下。
\end{itemize}

采用统一界面后，读写文件的动作可以操作几乎所有资源，程序不必为“读磁盘”“读传感器”“读另一进程的状态”分别实现接口。

每个运行中的程序默认打开三个通道。

\begin{center}
\small
\begin{tabularx}{\textwidth}{L{2.4cm} L{1.4cm} L{2.4cm} X}
\toprule
名称 & 编号 & 默认去向 & 用途 \\
\midrule
标准输入 & 0 & 键盘 & 程序从这里读数据 \\
标准输出 & 1 & 终端 & 程序的正常结果 \\
标准错误 & 2 & 终端 & 错误与警告 \\
\bottomrule
\end{tabularx}
\end{center}

重定向用于改变这三个通道的去向。

\begin{shellsession}
$ ls > files.txt          # 标准输出写入文件，屏幕不显示
$ echo note >> files.txt  # 追加，不覆盖
$ make 2> build.err       # 只把错误信息写入文件
$ ./run.sh > out.log 2>&1 # 正常输出与错误写入同一文件
$ wc -l < files.txt       # 输入来自文件，不再来自键盘
\end{shellsession}

管道只连接标准输出。左侧程序报错时，错误信息仍直接出现在终端，不进入右侧程序。排查问题时可据此区分“没有结果”与“执行失败”。

\section{退出码与静默输出}

不输出是 Unix 工具的通用行为：命令成功时通常不打印任何内容。

\begin{shellsession}
$ mkdir demo
$
\end{shellsession}

这是为了保证可组合性。若每个命令成功时都输出“已创建”“已完成”，管道中会充满与结果无关的文本。

命令是否成功由退出码表示。

\begin{shellsession}
$ mkdir demo; echo $?
0
$ mkdir /root/nope; echo $?
mkdir: cannot create directory '/root/nope': Permission denied
1
\end{shellsession}

0 表示成功，非 0 表示失败；具体数值由各程序定义，“非零即失败”是通用约定。shell 可以依据退出码决定是否继续执行。

\begin{shellsession}
$ ./build.sh && ./deploy.sh   # 前一条成功才执行后一条
$ ./check.sh || echo "检查未通过"
\end{shellsession}

脚本依赖这一约定判断流程是否中断。第十章说明 \texttt{set -euo pipefail} 的作用：使脚本在任一步失败时停止，而不是带着错误继续执行。

\section{固定约定}

Unix 程序通常不需要配置文件即可运行，默认值覆盖多数场景。与之配套的是一组固定约定。

\begin{itemize}
  \item \textbf{帮助开关}：几乎每个命令都接受 \texttt{-{}-help}，多数接受 \texttt{-h}。
  \item \textbf{长短选项}：\texttt{-v} 与 \texttt{-{}-verbose} 通常等价，短选项用于输入，长选项用于阅读。
  \item \textbf{版本}：\texttt{-{}-version}。
  \item \textbf{试运行}：\texttt{-{}-dry-run}，报告将要执行的操作而不实际执行。危险操作前应先运行一次。
  \item \textbf{用户配置}：位于 \texttt{\textasciitilde/.config/} 下，一个程序一个目录。
  \item \textbf{系统配置}：位于 \texttt{/etc/} 下，文件名通常与程序名相同。
  \item \textbf{日志}：由系统统一管理，不由各程序自行保存在私有目录。
\end{itemize}

\texttt{\textasciitilde} 表示家目录。\texttt{/etc} 存放系统级配置，\texttt{/var} 存放变化的数据，\texttt{/usr} 存放程序与只读资源。这套布局称为文件系统层次标准（FHS），第五章详述。

\begin{culturenote}
这些约定的作用是让使用者能够推测。面对一个未见过的命令，可以预期它有 \texttt{-{}-help}、支持 \texttt{-{}-dry-run}、配置文件位于 \texttt{/etc} 或 \texttt{\textasciitilde/.config}。可推测意味着不必每次查阅文档。
\end{culturenote}

\section{手册页}

\texttt{man} 是随系统分发的手册，不依赖网络与浏览器，服务器上必然存在。

\begin{shellsession}
$ man tar          # tar 的完整手册
$ man 5 passwd     # 第 5 节：配置文件格式
$ man -k compress  # 按关键词搜索全部手册
$ tldr tar         # 只看常用示例（需另行安装）
\end{shellsession}

手册分为若干节，同一名称可能出现在多节中。第 1 节是命令，第 5 节是文件格式，第 8 节是系统管理命令。\texttt{man 5 passwd} 与 \texttt{man passwd} 的内容不同。

进入手册后，\texttt{/关键词} 搜索，\texttt{n} 跳至下一处匹配，\texttt{q} 退出。

手册的常规结构为：名称、用法概要、描述、选项、示例、相关文件、已知问题。有效的读法是从示例入手，再回查选项，不按顺序通读。

\begin{lab}[练习：用三种方式查同一件事]
\begin{enumerate}
  \item 运行 \texttt{tar -{}-help}，找到“解压到指定目录”对应的选项。
  \item 运行 \texttt{man tar}，用 \texttt{/extract} 搜索，找到同一选项的完整说明。
  \item 运行 \texttt{man -k "disk usage"}，列出与本机磁盘占用相关的命令。
\end{enumerate}
三种方式回答同一问题，详细程度不同。日常先用 \texttt{-{}-help}，不足时再查 \texttt{man}。
\end{lab}

\section{提问与缺陷报告}

Linux 的工具由社区维护。使用中会遇到缺陷，也需要他人协助。社区沿用的规则可以概括为一条：让他人能够复现问题。

一份可用的缺陷报告包含五项内容。

\begin{enumerate}
  \item \textbf{环境}：发行版与版本、内核版本、软件版本。NixOS 上还需给出相关配置片段。
  \item \textbf{步骤}：从干净状态开始，能够逐步重放的命令。
  \item \textbf{预期}：按文档应当出现的结果。
  \item \textbf{实际}：实际出现的结果，包含完整报错。
  \item \textbf{证据}：日志、截图、\texttt{-{}-version} 的输出。
\end{enumerate}

提问时容易被省略的一项是“已经尝试过的排查”。它的作用是避免回答者重复已被排除的方向。这部分内容不承担其他作用。

\begin{culturenote}
社区中的“RTFM”（去读手册）常被视为不友好，但它指向一个事实：多数问题在手册、\texttt{-{}-help} 与日志中已有答案。先查证再提问，节省的是双方的时间。这条规则在 AI 时代同样适用：Agent 也依赖文档，提供的上下文越准确，得到的回答越可靠。
\end{culturenote}

\section{时间线}

这些约定在几十年中逐层累积形成。

\begin{figure}[htbp]
\centering
\begin{tikzpicture}[font=\footnotesize, x=1mm, y=1mm]
  \draw[plum!55, line width=2.4pt] (12,0) -- (34,0);
  \draw[teal!55, line width=2.4pt] (34,0) -- (82,0);
  \draw[ochre!65, line width=2.4pt] (82,0) -- (122,0);
  \draw[sage!65, line width=2.4pt] (122,0) -- (148,0);
  \foreach \x/\year/\era in {%
    12/1969/plum, 34/1973/plum, 61/1983/teal, 82/1991/teal,
    100/1993/ochre, 122/2003/ochre, 146/今天/sage} {
    \draw[ink, line width=0.8pt] (\x,-1.6) -- (\x,1.4);
    \fill[\era] (\x,0) circle (1.3mm);
    \node[above=1.8mm, text=ink, font=\bfseries\scriptsize] at (\x,1.6) {\year};
  }
  \node[below=2mm, align=center, text width=20mm] at (12,-1.6) {Unix\\ 贝尔实验室};
  \node[below=2mm, align=center, text width=20mm] at (34,-1.6) {用 C 重写\\ 可移植};
  \node[below=2mm, align=center, text width=20mm] at (61,-1.6) {GNU 项目\\ 自由软件};
  \node[below=2mm, align=center, text width=20mm] at (82,-1.6) {Linux\\ 内核};
  \node[below=2mm, align=center, text width=20mm] at (100,-1.6) {Debian\\ 社区发行版};
  \node[below=2mm, align=center, text width=20mm] at (122,-1.6) {Nix\\ 可复现部署};
  \node[below=2mm, align=center, text width=20mm] at (146,-1.6) {NixOS};
  \foreach \i/\c/\t in {0/plum/起源, 1/teal/自由软件与内核, 2/ochre/发行版, 3/sage/可复现} {
    \fill[\c] (30+\i*32,-9) rectangle ++(3,-3);
    \node[right=1mm, font=\scriptsize, text=muted] at (33+\i*32,-10.5) {\t};
  }
\end{tikzpicture}
\caption{从 Unix 到 NixOS。四种颜色对应四个阶段，每一代都在解决同一类问题：软件的安装、升级、回滚与环境复现}
\end{figure}

1969 年，贝尔实验室的肯·汤普森与丹尼斯·里奇编写了一个小型操作系统，取名 Unix。它的三项决定影响了后续所有系统：以文件为统一接口、以小工具组合完成任务、以纯文本作为程序间的数据格式。

1973 年，Unix 用 C 语言重写。操作系统因此可以移植到不同硬件，系统代码也成为普通人可读、可改的文本。

1983 年，理查德·斯托曼发起 GNU 项目，目标是编写一套自由的操作系统。今天使用的 \texttt{ls}、\texttt{grep}、\texttt{sed}、\texttt{bash} 多数来自 GNU。

1991 年，林纳斯·托瓦兹发布 Linux 内核。GNU 工具与 Linux 内核结合，形成通常所说的 Linux 系统。

1993 年，Debian 项目启动，把“由社区长期维护的发行版”变成可持续的工程。

2003 年，Nix 的论文发表，提出用函数式方法描述软件部署：每个包拥有唯一的哈希路径，环境由表达式计算得出，不依赖安装顺序。NixOS 是这一思路的完整实现。

\begin{culturenote}
这条时间线上的每一代都在处理同一组问题：如何安装、如何升级、如何回滚、如何让他人复现相同环境。NixOS 不是例外，它是在“可复现”这一目标上走得最远的一次尝试。
\end{culturenote}

\section{旧工具与新工具}

\begin{center}
\small
\begin{tabularx}{\textwidth}{L{3.4cm} L{2.6cm} X}
\toprule
旧工具 & 新工具 & 差异 \\
\midrule
\texttt{grep -r} & \texttt{rg} & 默认递归，默认忽略无关文件，默认输出行号与颜色 \\
\texttt{find} & \texttt{fd} & 语法更短，默认忽略隐藏目录与 \texttt{.gitignore} 中的路径 \\
\texttt{cat} & \texttt{bat} & 语法高亮与行号，长文件自动分页 \\
\texttt{ls} & \texttt{eza} & 颜色、图标、树状视图、Git 状态 \\
\texttt{sed -i} & \texttt{sd} & 正则语法接近主流写法，替换前可预览 \\
\texttt{du -sh} & \texttt{dust} & 以条形图显示目录占用 \\
\texttt{top} & \texttt{btop} & 彩色界面，支持鼠标，指标更完整 \\
\texttt{man} & \texttt{tldr} & 只列常用示例 \\
\texttt{diff} & \texttt{delta} & 并排显示与语法高亮，可配合 Git 使用 \\
\bottomrule
\end{tabularx}
\end{center}

\begin{toolswap}{grep -rn "TODO" .}{rg -n TODO}
\texttt{rg} 不需要 \texttt{-r}，默认递归；不需要排除 \texttt{node\_modules}，默认遵循 \texttt{.gitignore}；文件名与行号是默认输出。
\end{toolswap}

这些工具都不是必需组件。合理的使用顺序是：先用旧工具完成任务，当同一环节第三次造成麻烦时，再替换为新工具。顺序颠倒会导致时间花在配置工具上，而不是解决问题。

\section{练习}

\begin{enumerate}
  \item 用管道统计当前目录下出现次数最多的三种文件扩展名。可用 \texttt{ls}、\texttt{awk}、\texttt{sort}、\texttt{uniq}。
  \item 找出本机占用空间最大的三个目录。先用 \texttt{du}，若已安装 \texttt{dust} 再用它比对结果。
  \item 对一个不存在的文件执行 \texttt{cat}，记录屏幕输出与 \texttt{echo \$?} 的结果。
  \item 打开 \texttt{man bash}，用 \texttt{/variable} 搜索到“变量赋值”一节后退出。目标是练习在手册中定位，不必通读。
  \item 为该次 \texttt{cat} 失败写一份缺陷报告草稿，要求他人按步骤可复现。
\end{enumerate}

\begin{checkpoint}
本章结束时，应能：
\begin{itemize}
  \item 说明一条管道中每个程序承担的步骤；
  \item 解释标准输入、标准输出、标准错误，并用重定向改变其去向；
  \item 依据退出码判断命令成功与否；
  \item 用 \texttt{-{}-help} 与 \texttt{man} 自行查到答案；
  \item 写出一份他人可复现的缺陷报告；
  \item 指出 Unix 的哪几项设计决定了今日命令行工具的形态。
\end{itemize}
\end{checkpoint}
